アンサンブル学習は、複数の機械学習モデルの予測を組み合わせることで、より高精度な予測を実現する手法です。
シンプルでありながら多くの場合に性能を向上させるため、様々な分野で広く用いられてきました。
また研究も盛んに行われてきており、様々な手法が提案されてきました。

しかしながら、アンサンブル学習には根本的な問いが残されています。
それは、\textbf{「アンサンブルの性能を決定づける要素とは何か」}という問いです。
この問いに答えることは、より良い手法を設計するために不可欠です。

私たちは、この謎を解明することに成功しました。
私たちはまず、情報理論の「Fanoの不等式」に着目しました。
Fanoの不等式は、雑音のある情報を復元する際の、誤り率の下限を与えるものです。
アンサンブルの文脈に置き換えると、アンサンブル手法の誤り率の下限(性能のバロメーター)を示すことになります。

私たちは、この下限が、３つの要素に分解されることを示しました:
\begin{enumerate}
\item 個々のモデルの精度。
\item モデルの間の多様性。
\item モデルの予測を組み合わせる際に失われる情報量(結合損失)。
\end{enumerate}
個々のモデルが高精度であることはもちろん重要です。
また、モデル間に多様性があれば、あるモデルの誤りを別のモデルが補うことができます。
しかし、複数の予測を組み合わせる際に、正しい予測が他の誤った予測に埋もれてしまうこともあります。
このような場合は情報の損失が大きくなり、アンサンブルの性能は十分に発揮されません。

以上のように私たちは、\textbf{アンサンブルの性能が精度・多様性・結合損失という3つの要素によって決まる}ことを解明しました。
私たちはさらに、さまざまなアンサンブル手法についてこれらの要素を計算し、どの手法がどの要素を改善するのかを分析しています。

なお余談ですが、この研究は、SemEval・CoNLLという自然言語処理の国際コンペティションで、アンサンブル手法を用いて好成績を収めた経験をきっかけに、始まったものです。
